\section{Web alkalmazások kezdetektől napjainkig}
A világháló (WWW) használatának kezdetét az 1989--1991 közötti periódusra datálják. Kezdeti úttörője Tim Berners-Lee volt, a CERN egyik munkatársa, aki a web alapvető komponensei közé tartozó HTTP-t (HyperText Transfer Protocol), HTML-t (HyperText Markup Language), URL-t, valamint az első webböngészőt és webszervert is megvalósította \cite{cern2024www35}. Az ekkor készült weboldalak főként statikusak voltak, majd a szerveroldali programozás és az adatbázisok használata lehetővé tette a dinamikus tartalom előállítását \cite{mdnServerSideIntro}. A 2000-es évektől a kliensoldali interaktivitás az AJAX és a böngészőoldali JavaScript használatával vált egyre hangsúlyosabbá \cite{garrett2005ajax}. A mobileszközök elterjedése a reszponzív dizájnt tette alapvető megközelítéssé \cite{mdnResponsiveDesign}, a skálázhatóságot pedig a felhőalapú számítástechnika segítette \cite{mell2011nistCloud}.

A modern fejlesztési folyamatok manapság nagy hangsúlyt fektetnek a DevOps szemléletre, a folyamatos integrációra, szállításra és telepítésre (CI/CD), az automatizált tesztelésre, valamint a konténerizációra \cite{dora2018accelerate}. Architekturális oldalon a felhőalapú és mikroszolgáltatás-alapú megoldások terjedése vált meghatározóvá \cite{mell2011nistCloud,fowlerLewis2014microservices}.
